\documentclass[11pt]{article}
\usepackage[margin=1in]{geometry}
\usepackage{amsmath,amssymb,amsthm}
\usepackage{times}
\usepackage{hyperref}
\usepackage{booktabs}
\usepackage{enumitem}
\usepackage{natbib}
\hypersetup{colorlinks=true,linkcolor=blue,citecolor=blue,urlcolor=blue}

\title{Ensemble Before Selection \\ \large Redundancy, Random Sampling, and the Limits of Recoverability}
\author{Flyxion \\ \small Independent Researcher}
\date{September 2026}

\begin{document}
\maketitle

\begin{abstract}
A gas at equilibrium permits its temperature and pressure to be estimated from a sufficiently large representative sample of microscopic degrees of freedom, without a demon that first sorts fast from slow, useful from useless. A recent empirical result in machine learning, Random Attention, reports the same structural fact in an unrelated substrate: a language model's key-value cache can be evicted uniformly at random, discarding the majority of its stored reasoning trace, without loss of downstream accuracy, provided the original prompt is protected. A third case, drawn from a large language model for RNA, extends the principle further: a population that is not live at query time, but was consumed once during pretraining and compressed into a model's parameters, can still supply a bounded, measurable fraction of the same kind of recoverability. All three results are instances of a single principle, one older than any of these fields and rarely stated in general form: random sampling reliably estimates a target quantity when that quantity is distributed redundantly across many sufficiently exchangeable carriers, whether those carriers are live or amortized into a learned representation, and it becomes unreliable, in the limit unable to reconstruct the quantity from the surrounding population at all, once recovery depends on preserving one particular, non-interchangeable carrier. This essay develops that principle, using thermodynamic ensembles as the founding case, the passcode experiment inside Random Attention as the decisive counterexample marking its boundary, and single-sequence coevolutionary inference in a pretrained RNA model as the case that reveals a third, intermediate regime. The result is a general criterion for determining when selection mechanisms, understood broadly as importance scorers or other procedures that rank candidates before discarding most of them, are doing genuine epistemic work, and when they are elaborate machinery for a task that uniform sampling, live or amortized, already solves.
\end{abstract}

\section{Introduction}

Two questions are easy to conflate and worth separating from the outset. The first is whether a system's macroscopic behavior can be predicted from a sample of its microscopic constituents. The second is whether that sample must be chosen carefully, by some scoring rule that identifies which constituents matter, or whether an arbitrary subset will do. Statistical mechanics answered the second question for physical ensembles more than a century ago: equilibrium macroscopic observables can be estimated from a sufficiently large representative sample of microscopic degrees of freedom, subject to standard ergodicity conditions, with sampling error contracting as the sampled population grows. No molecule needs to be flagged in advance as important. Although molecules occupy distinct microstates, none has a privileged identity relative to the macroscopic observable being estimated. The macroscopic observable is recoverable because it is not, in any meaningful sense, a property of individual particles at all; it is a property of the distribution, and the distribution is what a sufficiently large random subsample already approximates.

A result published in 2026 on key-value cache management in large language models reports what looks, on first inspection, like an unrelated finding in an unrelated field. During long chains of reasoning, a transformer's key-value cache grows linearly with the number of generated tokens, and serving the model under a fixed memory budget requires discarding, or evicting, most of that cache as generation proceeds. Existing eviction methods have followed a single paradigm: assign every cached token a score estimating its future importance, and keep the highest-scoring entries. \citet{wang2026random} test this paradigm directly by replacing the scoring step with nothing at all. Random Attention protects the original prompt in full and evicts every other cached position uniformly at random, independently in each attention head, computing no score whatsoever. Across four models and six reasoning benchmarks in mathematics, science, and code, this null policy matches the strongest scored evictor while running 32 to 43 percent faster in deployment, because it never performs the scoring pass those methods require.

The claim of this essay is that these two results instantiate the same abstract pattern. Their physical implementations and probabilistic structures differ substantially: molecular exchangeability under equilibrium conditions is a much stronger and better-characterized symmetry than the contingent textual and cross-head redundancy observed in current reasoning models. What the two share is not a common probability model but a common structural condition. Both concern populations of carriers, molecular degrees of freedom in one case, and cached positions represented differently across attention heads in the other, from which a target quantity (pressure; the information needed to continue reasoning correctly) can be recovered without first identifying one uniquely important carrier, because the quantity in question is not concentrated in any particular carrier. Its evidential support is distributed across the population rather than concentrated in one indispensable record, and a uniform sample inherits that distributed support automatically.

The essay's second claim is that this principle has a sharp and experimentally demonstrated boundary, and the boundary is at least as important as the principle itself. Random Attention's own authors run a control in which a single fact, a synthetic passcode, is stated exactly once deep inside a reasoning trace and never restated, then queried much later after many rounds of eviction. Under the tested budget and protocol, Random Attention retrieves the passcode in none of the reported trials, whereas a history-sensitive scoring method recovers it approximately 84 percent of the time. This is the tagged particle inside the gas: the single molecule whose individual trajectory, rather than the ensemble's statistics, is the observable of interest. No amount of uniform sampling reconstructs a quantity that was never made redundant across the sampled population in the first place.

The result is a general criterion, stated in Section 6, for distinguishing quantities that are ensemble-recoverable, whether the population supplying that recoverability is live or amortized into a learned representation, from quantities that require identity-preserving selection, together with an account of what a selection mechanism (Maxwell's demon, an attention score, any importance-ranking procedure) is actually purchasing when it earns its computational cost.

\section{Heat as Ensemble Recoverability}

Temperature and pressure are defined, in kinetic theory, as statistics of a distribution over molecular velocities and positions, not as properties any individual molecule possesses. A single molecule has a velocity; it does not have a temperature. Temperature is what emerges when the velocities of very many molecules are averaged, in a particular way, over the ensemble. This has a direct consequence for measurement. To estimate a gas's pressure, an instrument does not need to interrogate every molecule, still less identify in advance which molecules are the important ones. It needs only a sufficiently large random sample, because the law of large numbers guarantees that the sample statistic converges to the population statistic regardless of which molecules happen to be included. Exchangeability is doing the entire work: any molecule can stand in for any other with respect to the quantity being measured, because the quantity being measured was never a fact about molecular identity to begin with.

This is the structural condition under which uniform sampling is an adequate null procedure: no prior ranking by carrier identity is required for the sample statistic to converge to the ensemble statistic. This is a weaker claim than optimality; stratified, importance-weighted, or variance-reducing sampling schemes can outperform a uniform draw even when the target is an ordinary ensemble statistic, and nothing here rules that out. What uniform sampling establishes is only that ranking is not \emph{necessary}: no molecule is indispensable by identity to the pressure estimate. Particular finite samples can fluctuate, and omitting an extreme observation can move their value, but the estimator does not require one antecedently designated molecule to remain present. Maxwell's thought experiment about a demon who sorts fast molecules from slow ones is illuminating here precisely because it describes a task that is possible in principle (a demon with perfect information about every molecule's velocity could, by selective sorting, do work that appears to violate the second law) and is nonetheless irrelevant to the ordinary business of measuring a gas's bulk properties, which requires no sorting at all. The demon's labor buys something (a temperature gradient, in Maxwell's version, or a targeted individual fact, in the argument below), but it does not buy a better estimate of the ensemble average, because the ensemble average was already available for free.

\section{Random Attention as a Second Instance}

The mechanics of \citet{wang2026random} translate this condition into a completely different physical substrate. A reasoning trace generated by a large language model is not, in general, a sequence in which each token contributes a unique, non-repeated piece of information. The paper's central diagnostic observation is that reasoning traces are redundant at two independent levels simultaneously.

The first redundancy is textual: a model working through a long derivation tends to restate its active premises and intermediate results as it proceeds, because each new step typically re-invokes the quantities it depends on. A value computed early in a derivation and needed later is, empirically, not stored once and referenced; it is recomputed, restated, or paraphrased at each point where it is used. The second redundancy is architectural and specific to the transformer mechanism: prior to eviction, every attention head maintains its own key-value representation of every cached position, and each head decides, independently of the others, which cached positions to discard. These per-head representations are not identical copies of one another; each head projects the same token into a different subspace, so a head's representation is better described as a distinct, partially overlapping encoding of the underlying information than as a duplicate. What is architecturally redundant is the number of independent encodings available, not the content of any two encodings being the same.

Given both redundancies, the paper's planted-fact probe shows directly that recovery of a value can depend on partial evidence pooling across several particular heads rather than on any one head holding a complete, retrievable copy. A fact held in exactly one attention head is retrieved in as few as one to three percent of trials; specific two-head combinations vary considerably in effectiveness (roughly sixteen, thirty-nine, and sixty percent depending on which pair), while selected three-head and all-head configurations rise to roughly eighty-three and ninety-nine percent, a pattern the paper describes as strongly superadditive. This establishes distributed, complementary preservation across heads, not literal replication of an equivalent copy in every head: some heads are considerably more informative witnesses than others, and the redundancy is one of sufficient substitutability across a population of non-identical encodings, not strict interchangeability. The disanalogy with the gas is worth keeping precise for this reason: molecules in an equilibrium ensemble are treated as exchangeable with respect to a macro-observable almost by definition, whereas attention heads are specialized, and the paper's contribution is to show they nonetheless provide enough distributed, complementary coverage for uniform per-head eviction to recover what the model needs.

This explains, without appeal to any special property of randomness itself, why Random Attention's uniform per-head eviction matches learned scoring rules on the reasoning trace once the prompt is separately protected. The scoring rules are not extracting information that a uniform draw misses; they are, at best, redundantly confirming a redundancy the trace already provides. The paper's own controlled experiment makes this explicit: removing cross-head diversity entirely, by forcing every head to share one randomly drawn keep-set instead of drawing independently, costs essentially nothing on real traces (0.871 versus 0.874 accuracy at one budget the authors test), because the textual redundancy alone already carries a restated copy of whatever the model still needs. The architectural redundancy across heads is not load-bearing when the textual redundancy is already sufficient; it becomes load-bearing only when the textual redundancy is absent; which is exactly the condition the passcode experiment is built to isolate.

The paper's finding about the prompt sharpens the same point from a different angle. The original prompt is externally supplied rather than generated as part of the model's self-refreshing working trace. Although portions of it may later be repeated during reasoning, its constraints are not guaranteed to be redundantly reinstantiated the way the model's own working state is; reconstruction from the surrounding trace is possible but not assured. It therefore constitutes the fragile portion of the cache the paper identifies, and the paper accordingly finds that whether a method protects the prompt by rule, rather than by leaving it to a score, accounts for most of the reported difference between competing scored evictors. Once every method is given the same prompt-protection rule, the residual gap between them, the paper reports, becomes small and in several cases favors the policy that ranks nothing over the policies that spend computation ranking. The prompt, in other words, sits closer to the identity-recoverable side of the boundary this essay is describing, while the reasoning trace built on top of it sits, for contemporary models, mostly on the ensemble-recoverable side. Distinguishing the two is not a detail of the paper's method; it is the paper's central structural discovery, restated here in the vocabulary of ensembles and identity.

\section{The Demon and the Passcode}

Maxwell's demon is worth taking seriously as more than a metaphor here, because the paper supplies a demon that actually earns its cost and reports exactly what it buys, with an important qualification attached. Among the scored eviction methods the paper compares against, one, called R-KV, computes an importance score by accumulating each cached position's attention weight over the entire history rather than over a recent window. This is architecturally the demon's operation, and it is worth being precise about what kind of operation it is: R-KV is not a smarter sampler, drawing a better-chosen random subset. It is a targeting mechanism, one that decides in advance, on content-sensitive grounds, which specific carrier to protect from eviction regardless of how long ago it appeared. The relevant contrast is not between two ways of sampling but between sampling at all and intervening to retain one designated record.

The experiment designed to isolate this capability states a synthetic passcode exactly once, embeds it fifty-seven eviction rounds before a query that requires it, and reports, for every method under comparison, both the fraction of trials in which the correct passcode is reproduced and the log-probability the model assigns to the correct value at the point of query. Random Attention, drawing uniformly with no accumulated history, reproduces the passcode in none of the reported trials, and the model's assigned log-probability to the correct value falls to a level the paper characterizes as equivalent to the fact being entirely absent from the cache. This is an empirical zero at the paper's tested sample size and retention budget, not a claim that recovery is mathematically impossible in every configuration; under a fixed keep-rate, a uniform draw does carry a nonzero, budget-determined probability of retaining any single position, including the passcode's, and that probability rises as the retained fraction grows. What a uniform draw cannot do, at any retained fraction short of the whole cache, is target the passcode specifically; its survival is incidental to the draw rather than a consequence of the mechanism recognizing what the passcode is. The recency-weighted scored methods, which weight recent attention far more heavily than distant attention, fare almost as badly as pure chance, at 0.4 to 1.6 percent retrieval, despite two of them being the strongest performers in the paper's aggregate benchmark comparisons. Only R-KV, whose score is specifically built to accumulate importance across the whole history rather than a recent window, recovers the passcode with any reliability, at 84 percent.

This is the tagged molecule, with the same qualification carried over. A gas's pressure is insensitive to which molecule carries which velocity, but if the observable of interest were instead "where is this specific labeled molecule right now," no amount of uniform ensemble sampling reconstructs the answer from the surrounding molecules, because the question is no longer about a distribution, it is about an individual, non-interchangeable carrier's continued identity; enlarging the sample only raises the odds that the labeled molecule happens to be swept up in it, which is brute-force protection rather than reconstruction from other evidence. The passcode is exactly this kind of observable inside a reasoning trace: a fact stated once, at one position, with no other carrier holding an equivalent copy, so that its recoverability depends entirely on whether that single position, or a targeting mechanism built to protect it, survives. A striking further observation the paper reports is that this needle-finding capability and aggregate benchmark strength are nearly orthogonal: R-KV, the strongest needle-finder in the passcode experiment, is described in the paper as leading only one column of its main results table, while TriAttention, the strongest aggregate performer among the scored baselines, recovers almost nothing in the passcode experiment. Nor should R-KV's success here be read as evidence that accumulated-attention scoring identifies semantic importance in general; the experiment shows only that this particular statistic, tracking sustained attention over the full history, happens to correlate with the survival of a planted, once-stated fact in this regime. A demon that is very good at tracking one labeled particle is not thereby a demon that improves the ensemble reading, and a scorer that is excellent at the ensemble task is not thereby equipped to track a single unrepeated fact. These are different jobs, answered by different mechanisms, and a benchmark that rewards only the aggregate task will systematically undervalue the second capability, precisely because real reasoning traces, as generated by current models, rarely pose the second kind of problem. The paper's own framing supports this: the case is rare on real traces because the model keeps restating what it is still using, so the demon's specialized skill goes largely unexercised outside of a deliberately constructed control.

\section{A Third Regime: Amortized Ensemble Access}

The two regimes developed so far, live ensemble access and singular identity, do not exhaust the space. A 2026 result on RNA sequence modelling identifies an intermediate case worth naming in its own right: a target quantity that is genuinely ensemble-recoverable in principle, but only from a population the system does not have live access to at the moment of use.

Classical coevolution analysis recovers structural contacts in a biomolecule from a multiple sequence alignment (MSA): many related sequences, treated as a population, whose correlated substitutions at paired positions reveal which residues interact physically. This is live ensemble recovery in the sense already developed here: the signal is distributed across many concurrently available, mutually substitutable sequences, and mutual information or direct-coupling analysis extracts it by pooling across that population at analysis time.

\citet{upadhyay2026nucleicbert} report a large language model for RNA, NucleicBERT, trained by masked-token prediction on approximately 30 million non-coding RNA sequences with no alignment information supplied during training or at inference. Once pretrained, the model can be probed with a masked-likelihood-influence (MLI) score: for a single input sequence, MLI measures how much masking one position degrades the model's confidence in the nucleotide at another position, yielding a coupling matrix computed entirely from that one sequence. Averaged post hoc over a family's aligned sequences purely to permit comparison, MLI recovers a measurable fraction of the same coevolutionary signal that mutual information and direct-coupling analysis extract from the full alignment: across fourteen Rfam families, MLI sits well above a random null on enrichment factor, but consistently below the MSA-grounded references it is compared against. The per-sequence signal is reported as roughly flat across each query's similarity to the pretraining corpus, which the authors offer as evidence against memorization: the model does not appear to be recalling a stored answer for sequences resembling its training data, but expressing something learned more generally about the family's constraint structure.

What has happened structurally is that the population needed for ensemble recovery, the many related sequences an alignment would supply, was not discarded but relocated. It was consulted once, during pretraining on millions of sequences, and compressed into the model's parameters; at inference time a single query sequence draws on that compressed population rather than on any population present in its own input. Call this amortized ensemble recovery: a target quantity that is population-recoverable in principle, accessed at query time through a single carrier whose capacity to stand in for the population was purchased in advance by prior exposure to many others.

Amortized recovery explains a pattern that recurs across NucleicBERT's downstream tasks and gives the distinction a further empirical test. Where a task supplies little labelled data of its own, such as secondary-structure or contact-map prediction, a frozen pretrained backbone with only a linear head attached already substantially outperforms a frozen, randomly initialized backbone of identical architecture, and the residual gain from full fine-tuning is comparatively modest; the amortized population access is reported as already doing most of the work before any task-specific gradient is taken. Where a task instead supplies abundant labelled data of its own, such as splice-site classification with tens of thousands of labelled examples per species, the gap between the frozen pretrained backbone and full fine-tuning widens sharply, and even a randomly initialized backbone reaches respectable accuracy once fine-tuned end to end: here the task's own supervision, rather than the amortized population, is doing most of the work, because enough live, task-specific signal is available to make the amortized shortcut largely unnecessary. Amortization matters, in other words, specifically where a live substitute for it, whether an alignment or a labelled dataset, is scarce.

The boundary of amortized recovery is measurable rather than assumed. Reported family by family, MLI's advantage over the random null is real but bounded, and it does not reach the ceiling set by an actual alignment: pretraining approximates the population it was shown, it does not reconstitute it on demand. This gives amortized recovery a signature distinct from either of the other two regimes developed so far. Unlike live ensemble recovery, its accuracy at a given query cannot be improved simply by supplying more sequences alongside that query, since no additional population is present at inference time to sample from; the population that mattered was already spent during pretraining. Unlike identity recovery, it is not reduced to chance either, since the compressed population continues to supply real, if attenuated, signal regardless of which single query is presented. It occupies the middle of the principle this essay has been developing, and it sharpens that principle rather than complicating it: what determines recoverability is not whether a population was ever consulted, but whether it is a population, live or amortized, that the current query can still draw on, as opposed to a fact that was never made redundant by any population, live or amortized, in the first place.

\section{Ensemble Recovery Versus Identity Recovery}

The three cases developed above, thermodynamic, computational, and single-sequence biological, are instances of a single general distinction, which can be stated without reference to any one substrate. A terminological note first: in statistical mechanics proper, an ensemble is the collection of possible microstates consistent with given macroscopic constraints, not the population of molecules in a single system. The argument of this essay is closer to a law-of-large-numbers claim about a sample drawn from a population than to an ensemble-average in that technical sense, and "population-recoverable" is the more exact term where precision matters. The word "ensemble" is retained in what follows because it names the shared structural feature cleanly, namely a target quantity's insensitivity to which particular carriers, molecules or heads or tokens or sequences, happen to constitute the sample; the technical content of the claim rests on exchangeability, not on the ensemble formalism as such.

\begin{align}
\text{live ensemble recovery} &: \text{a target quantity is estimable from many sufficiently exchangeable carriers, jointly present at query time}, \notag \\
\text{amortized ensemble recovery} &: \text{a target quantity is estimable from a single carrier at query time, via a mechanism fit in advance to many other carriers}, \notag \\
\text{identity recovery} &: \text{a target quantity is preserved by one designated, non-interchangeable carrier}. \notag
\end{align}

Where a quantity is population-recoverable, live or amortized, a selection mechanism, whether a scoring function over cached tokens or a demon sorting molecules by velocity, purchases nothing that uniform random sampling does not already supply as an unbiased default, because the quantity was never concentrated in the carriers the mechanism is choosing between. This is the condition under which Random Attention's null policy matches learned scores, it is the condition under which a sufficiently large random subsample of a gas yields an accurate pressure reading, and it is the condition, in its amortized form, under which a single RNA sequence yields a bounded but real coevolutionary signal without any alignment supplied at query time. Where a quantity is identity-recoverable, in contrast, uniform sampling does not permit reconstruction of the quantity from the remaining population, no matter how large the retained sample is drawn from carriers that do not individually hold it. Increasing the retained fraction does raise the probability that the one relevant carrier itself survives, to the retained fraction in a single round of uniform sampling without replacement, and that probability compounds downward under repeated rounds of eviction; but this is brute-force protection by enlarging what is kept, not extraction of the target from interchangeable evidence, and a sufficiently long or tightly budgeted elimination process drives it toward zero regardless of how the sample was drawn. A targeting mechanism, by contrast, can raise a single carrier's survival probability without enlarging the retained set at all. This is the condition under which the passcode was recovered in none of Random Attention's reported trials while a scorer built to track exactly that carrier recovered it reliably.

The amortized case occupies a distinct position in this scheme rather than a compromise between the other two. Its ceiling is set by the population it was fit to, which is why NucleicBERT's single-sequence coevolutionary signal, however real, does not reach the strength of mutual information or direct-coupling analysis computed on an actual alignment: amortization approximates a population, it does not reconstitute one on demand. Its floor is set by the fact that a compressed population, unlike an absent one, still supplies genuine signal at every query, which is why the amortized case never collapses to the passcode's zero. Between these two bounds, what determines how much amortization is worth, in the RNA model's own results, is how much live population access a task's own labelled data can supply as a substitute: amortization does heavy lifting where labels are scarce and comparatively little where they are abundant, because abundant labels are themselves a second route to the same underlying population information.

The general principle can be stated compactly:

\[
\boxed{\text{Random sampling estimates replicated structure, live or amortized; it preserves singular identity only by chance.}}
\]

A formal statement sharpens why the two extremal regimes, live ensemble recovery and identity recovery, differ in kind rather than only in degree. For an ensemble statistic estimated from $K$ sufficiently weakly correlated carriers, sampling error commonly contracts on the order of
\[
O\!\left(K^{-1/2}\right).
\]
For a singular carrier sampled uniformly without replacement from $N$ candidates, by contrast, one-round survival is merely
\[
\Pr(\text{survival}) = \frac{K}{N}.
\]
The first relation expresses reconstruction from distributed evidence: enlarging $K$ improves the estimate by aggregation, and the estimate remains recoverable even if any particular carrier is missing. The second relation expresses only the chance preservation of an irreplaceable record: enlarging $K$ raises the odds that the one relevant carrier happens to be retained, but supplies no mechanism by which its content could be reconstructed from the other $K-1$ carriers if it is not. The amortized case does not reduce cleanly to either formula, and that is itself informative: its error is bounded not by the size of a population present at query time, since none is present, but by the size and composition of the population it was fit to beforehand, a quantity fixed at pretraining and no longer adjustable by anything done at inference.

Three clarifications keep this claim from overreaching. First, the distinction is a property of the target quantity relative to a given population of carriers, not an intrinsic property of a domain. The same reasoning trace that is ensemble-recoverable with respect to "what does the model still need in order to continue correctly" is identity-recoverable with respect to "what exact passcode was stated at position 4,200," because these are different target quantities defined over the same underlying data. A gas is ensemble-recoverable with respect to pressure and identity-recoverable with respect to the trajectory of one tagged particle. A single RNA sequence is amortized-recoverable with respect to its family's structural contacts, given a model pretrained on a large corpus of related sequences, and would be closer to identity-recoverable with respect to the same contacts given no such pretrained model at all. The question "is selection necessary here" cannot be answered from the domain alone; it requires specifying what is being recovered, and by what, if anything, it is being recovered.

Second, redundancy is not a fixed property of a system but a contingent, and in both computational cases an engineered, one. The Random Attention paper is explicit that the textual redundancy it exploits is a property of how current reasoning models happen to generate chains of thought under standard training, not a logical necessity of reasoning as such. A model, or a training procedure, could in principle produce traces that state each intermediate fact exactly once, in which case the reasoning trace would shift from the ensemble-recoverable side of the boundary to the identity-recoverable side, and Random Attention's parity with scored methods would disappear along with it. The paper's own account of where its method loses ground, specifically on code-generation tasks with long, largely non-redundant prompts, illustrates this directly: prompts in that setting consume a large and variable share of the eviction budget precisely because they are not restated, and the paper notes that a smarter rule might compress such prompts rather than protect them whole, which is future work the null hypothesis itself cannot address. The amortized case is contingent in an analogous way, on the composition of the pretraining corpus rather than on a single trace's structure: the reported flatness of MLI quality across a query's similarity to the training data supports the claim that the model is not simply memorizing near-duplicates, but this is a claim about the fourteen evaluated families and the corpus actually used, not a guarantee that hold for a family unlike anything the model was shown. The boundary between the regimes is therefore an empirical fact about a given system's redundancy structure, live or amortized, at a given time, not a permanent classification of reasoning, cognition, or physical systems in general.

Third, the amortized case's dependence on prior population access has a further consequence for where selection remains valuable even under amortization. Nothing about compressing a population into pretrained parameters removes the need for a targeting mechanism on quantities that were never made redundant by that population in the first place; a passcode-like fact, planted once with no analogue anywhere in the pretraining corpus, would remain as unrecoverable to an amortized model as to a live ensemble with no access to it. Amortization changes what counts as available population, not the underlying logic of the boundary between population access and singular identity.

\section{Conclusion}

The apparent success of a selection mechanism, whether a physical demon, an attention-importance score, or any procedure that ranks candidates before discarding most of them, is not by itself evidence that selection was necessary. It is evidence only that the mechanism did not make things worse, which is a much weaker claim, and one satisfied automatically whenever the target quantity was already recoverable from an arbitrary subsample, live or amortized. The question worth asking of any such mechanism is not whether it works, since a sufficiently redundant target will make almost anything work, but what would have happened to a uniform random baseline given the same budget and the same protection of whatever carriers are individually irreplaceable. Random Attention is valuable, beyond its specific engineering contribution, for having asked and answered exactly that question, and for having supplied, in the planted passcode, a clean demonstration of where the answer flips. NucleicBERT's single-sequence coevolutionary signal extends the same question into a regime where the relevant population is no longer present at all at the moment of use, and shows that the answer need not be all-or-nothing there either: a compressed population can supply real, bounded recoverability, measurably short of what a live population would supply, and measurably above what an absent one would. Statistical mechanics supplied the founding mathematical case for the two extremal answers: representative sampling estimates the macrostate, whereas a designated microscopic history cannot be reconstructed from the macrostate that deliberately abstracts it away. That the computational cases hold for present-day reasoning models and present-day pretrained sequence models is itself a contingent fact about how those systems are currently trained; should reasoning traces grow more informationally economical, or should a pretraining corpus fail to cover a given family at all, the boundaries this essay has drawn would shift accordingly, and selection would recover the ground it presently cedes to chance or to amortization.

\bibliographystyle{plainnat}
\begin{thebibliography}{9}

\bibitem[Wang et al.(2026)]{wang2026random}
Wang, H., Qiu, J., Zhao, W., Qian, C., Yang, L., Han, J., Ji, H., Savarese, S., Heinecke, S., and Wang, H. (2026).
\textit{Random Attention: Rethinking KV Cache Eviction for Efficient Reasoning}.
Salesforce AI Research. arXiv:2609.03430.

\bibitem[Maxwell(1871)]{maxwell1871}
Maxwell, J. C. (1871). \textit{Theory of Heat}. Longmans, Green, and Co.

\bibitem[Cai et al.(2025)]{cai2025rkv}
Cai, Z., Xiao, W., Sun, H., et al. (2025).
R-KV: Redundancy-aware KV cache compression for reasoning models.
\textit{The Thirty-ninth Annual Conference on Neural Information Processing Systems}.

\bibitem[Upadhyay et al.(2026)]{upadhyay2026nucleicbert}
Upadhyay, U., Herold, J., G\"otz, M., and Schug, A. (2026).
NucleicBERT interprets RNA sequence space through self-supervised language modelling.
\textit{Nature Machine Intelligence}. https://doi.org/10.1038/s42256-026-01295-9.

\end{thebibliography}

\end{document}

